% TOP_PROBLEM: {"id": 30006985, "title": "Huisken's Generic Singularity Conjecture for Mean Curvature Flow", "category_id": 9, "category": {"id": 9, "name": "pde", "display_name": "Partial Differential Equations", "description": "PDEs and their applications in physics and geometry.", "slug": "pde", "order_index": 9, "created_at": "2026-07-31T15:26:25.671Z"}, "status": "open"}
% FIRST_POSTED: 2026-09-27
% ATTEMPT_STATUS: unresolved
% !TeX program = lualatex
\documentclass[11pt]{article}
\usepackage[margin=25mm]{geometry}
\usepackage{fontspec}
\setmainfont{Latin Modern Roman}
\usepackage{amsmath,amssymb,mathtools}
\usepackage{unicode-math}
\setmathfont{Latin Modern Math}
\usepackage{xurl,hyperref}
\hypersetup{colorlinks=true,urlcolor=blue}
\setcounter{secnumdepth}{0}
\setlength{\parindent}{0pt}
\setlength{\parskip}{5pt}
\setlength{\emergencystretch}{3em}
\begin{document}
\section*{\texorpdfstring{Huisken's Generic Singularity Conjecture for Mean Curvature Flow}{Huisken's Generic Singularity Conjecture for Mean Curvature Flow}}\label{30006985}
\subsection{Definitions and mathematical statement}
Mean curvature flow evolves an embedded hypersurface by $\partial_tF=-H\nu$, with $H$ its scalar mean curvature. A singularity model is a tangent flow obtained by parabolic rescaling around a singular point. The conjectural permitted models are multiplicity-one round shrinking spheres and generalized cylinders, whose time-$-1$ slices are
\[
 S^k(\sqrt{2k})\times{\mathbb{R}}^{n-k}\subset{\mathbb{R}}^{n+1},\qquad1\le k\le n.
\]
The catalogue requests this for every singularity arising from an open dense generic set of smooth closed embedded initial hypersurfaces. A precise global statement requires the source's topology on initial hypersurfaces and its weak-flow continuation convention after singularities. Classifying stable self-shrinkers alone does not establish that all singularities of every member of an open dense set have these models.
\par\smallskip\textit{Formulation and concept references:} \href{https://arxiv.org/abs/0908.3788}{[S1]}.

\subsection{Short English statement}
For an open dense set of smooth closed initial hypersurfaces, must every mean-curvature-flow singularity look like a single shrinking round sphere or cylinder?
\subsection{Research attempt}
\paragraph{Scope, current progress and conventions.}
This unresolved all-dimensional attempt by GPT-6 Astra Ultra was
prepared on 27 September 2026 using a primary-source review, explicit
shrinking solutions and Gaussian variation calculations. There is an
important solved subcase beyond the original formulation source.
Chodosh--Choi--Mantoulidis--Schulze [E1], Theorem 1.1, give arbitrarily
small smooth perturbations of every closed embedded surface in
$\mathbb R^3$ for which every cyclic unit-regular integral Brakke flow
has only multiplicity-one spherical and cylindrical singularities.
Their Remark 1(c) supplies openness as well as density. This resolves
the surface case with their continuation convention. It is not an
all-dimensional theorem. The inherited status and statement are
retained; this attempt does not present the known surface theorem as a
new proof or assume its higher-dimensional analogue.

For the direct computations take a smooth closed hypersurface before
its first singular time, with outward normal on a round sphere and
$H>0$ there. Thus the velocity is $-H\nu$. Assertions about tangent
flows and continuation after singularities require weak-flow theory;
the elementary computations below do not construct that theory.

\paragraph{Exact shrinking models and the normalization of their radii.}
On a round $n$-sphere of radius $R(t)$, the principal curvatures are
$1/R$, so
\[
             R'=-\frac nR,\qquad R(t)^2=R(0)^2-2nt.
\]
At its extinction time $T=R(0)^2/(2n)$, rescaling by
$(T-t)^{-1/2}$ gives the sphere of radius $\sqrt{2n}$. For the product
$S^k(R)\times\mathbb R^{n-k}$, exactly $k$ principal curvatures equal
$1/R$ and the others vanish. Hence $R'=-k/R$, and the analogous
rescaled radius is $\sqrt{2k}$. These computations reproduce the root
statement with its factor $2$, rather than a convention-dependent
radius $\sqrt{k}$.

A time-$-1$ shrinker obeys
\[
                       H=\frac{\langle x,\nu\rangle}{2}.       \tag{1}
\]
For the cylinder the right side is $R/2$, while the left is $k/R$;
(1) again gives $R^2=2k$. The flat case $k=0$ is a stationary plane,
which models a regular multiplicity-one point rather than one of the
genuine singularities requested in the statement. These explicit
flows produce the permitted models. They do not prove that another
initial hypersurface can only develop these models.

\paragraph{Deriving the Gaussian monotonicity identity with the chosen signs.}
Fix $(x_0,t_0)$ and write $r=x-x_0$, $\tau=t_0-t>0$. Define
\[
 \rho(x,t)=(4\pi\tau)^{-n/2}\exp\left(-\frac{|r|^2}{4\tau}\right),
 \qquad \Phi(t)=\int_{M_t}\rho\,d\mu_t.
\]
The area element evolves by $\partial_t d\mu=-H^2d\mu$.
Differentiating $\rho$ along the moving hypersurface gives
\[
 \frac{1}{\rho}\frac{d\rho}{dt}
 =\frac n{2\tau}-\frac{|r|^2}{4\tau^2}
       +\frac{H\langle r,\nu\rangle}{2\tau}.
\]
For the tangential vector $r^T$, direct differentiation yields
$\operatorname{div}_M r^T=n-H\langle r,\nu\rangle$ and
$\nabla_M\rho=-\rho r^T/(2\tau)$. Since $M_t$ is closed,
\[
 0=\int_{M_t}\operatorname{div}_M(\rho r^T)
  =\int_{M_t}\rho\left(n-H\langle r,\nu\rangle
                         -\frac{|r^T|^2}{2\tau}\right).
\]
Subtracting this zero identity with coefficient $1/(2\tau)$ from the
differentiated integral completes the square:
\[
 \Phi'(t)=-\int_{M_t}
       \left(H-\frac{\langle r,\nu\rangle}{2\tau}\right)^2
            \rho\,d\mu_t\leq0.                              \tag{2}
\]
This verifies the sign in (1) independently of a remembered convention.
For the centered shrinking sphere with $t_0=T$, the integrand vanishes
identically and $\Phi$ is constant. A cylinder satisfies the same
identity after the usual integrable Gaussian cutoff in its Euclidean
directions.

Equation (2) explains why a sufficiently regular rescaled limit with
constant Gaussian density is a self-shrinker. Establishing existence,
compactness, integrality and the appropriate equality case for arbitrary
tangent flows uses additional theorems. Even after those inputs, (1)
is a differential equation admitting a larger class of solutions than
the desired spheres and cylinders. Monotonicity is therefore a
reduction to shrinkers, not their classification or generic avoidance.

\paragraph{Gaussian mass detects multiplicity as well as shape.}
For a hypersurface $\Sigma$ put
\[
           F(\Sigma)=(4\pi)^{-n/2}
                        \int_\Sigma e^{-|x|^2/4}\,d\mu.
\]
On $\Sigma_k=S^k(\sqrt{2k})\times\mathbb R^{n-k}$, the flat Gaussian
integral cancels the factor from the Euclidean directions, leaving
\[
       F(\Sigma_k)=|S^k(1)|
              \left(\frac{k}{2\pi}\right)^{k/2}e^{-k/2}.
                                                               \tag{3}
\]
In particular $F(\Sigma_1)=\sqrt{2\pi/e}$ and
$F(\Sigma_2)=4/e$. A varifold consisting of the same support with
integer multiplicity $m$ has Gaussian mass $mF(\Sigma_k)$. Thus
recognizing the round support of a tangent object does not establish
multiplicity one. A proof excluding other smooth shrinker shapes
would still leave this separate obstruction. Small perturbations of
initial data cannot simply be presumed to remove multiplicity without
a theorem linking the perturbations to the weak tangent-flow measure.

\paragraph{A spherical second-variation calculation and a stability trap.}
For a smooth shrinker, the normal second variation at fixed Gaussian
center and scale has the form
\[
 F''(f\nu,f\nu)=-(4\pi)^{-n/2}\int_\Sigma fLf\,e^{-|x|^2/4},
 \quad L=\Delta_\Sigma-\tfrac12\langle x^T,\nabla_\Sigma\rangle
                  +|A|^2+\tfrac12.                           \tag{4}
\]
Here (4) is the standard Gaussian variational formula used in [S1]; we
compute its spectrum only on the round sphere. At radius $\sqrt{2n}$,
$x^T=0$ and $|A|^2=1/2$, so $L=\Delta+1$. The spherical harmonics of
degree $\ell$ have Laplace eigenvalue
$\ell(\ell+n-1)/(2n)$, and hence $L$-eigenvalue
\[
                   \mu_\ell=1-\frac{\ell(\ell+n-1)}{2n}.
                                                               \tag{5}
\]
The constant mode has $\mu_0=1$, the linear modes have $\mu_1=1/2$,
and every mode with $\ell\geq2$ has $\mu_\ell<0$. Thus the quadratic
form in (4) is positive on the subspace orthogonal to constants and
linear harmonics, while it is negative in the dilation and translation
normal directions. The scale test can also be seen directly: for a
round radius-$R$ sphere, $\log F$ differs from
$n\log R-R^2/4$ by a constant, whose second derivative at
$R=\sqrt{2n}$ equals $-1$.

In particular the statement ``the sphere minimizes the fixed-center
Gaussian functional under every variation'' is false. Generic
singularity stability must account for the freedom to recenter and
rescale. Entropy is the supremum of Gaussian functionals over those
parameters, not the single fixed functional in (4). The sphere
calculation is consistent with the correct treatment of symmetry
modes, but does not independently establish the global entropy
classification in [S1]. No spherical spectral calculation classifies
arbitrary noncompact shrinkers.

\paragraph{The missing step is global and involves quantifiers.}
Suppose for each unwanted shrinker one finds a nearby initial
perturbation that avoids that particular model near one spacetime
point. This has the form ``for each obstruction there exists a
perturbation.'' The target requires one perturbed initial hypersurface
whose entire continued flow avoids every unwanted model at every
singular point. Passing between these quantifiers needs compactness,
control of new singularities and a mechanism to globalize the local
avoidance. The new flow might develop an unwanted singularity elsewhere,
or a sequence of models might accumulate outside the region where a
chosen local estimate is uniform.

A decreasing Gaussian quantity alone does not force termination of a
perturbation procedure. Positive drops $2^{-j}$ have a finite sum, so
infinitely many strict decreases can fit within a bounded initial
budget. One needs a suitable uniform drop on controlled classes, a
finite stratification, or another proved argument ruling out an
infinite obstruction sequence. Such hypotheses must be supplied at
the dimensions, entropy bounds and regularity levels under discussion;
they cannot be imported from the surface case without checking them.

There is also a topology distinction. A countable intersection of open
dense sets can be dense without being open, as the irrational numbers
in $\mathbb R$ illustrate. Therefore a Baire-category construction of
flows avoiding successively many bad events, even if valid, does not
alone prove the root's open dense conclusion. Openness requires
stability of the good global singularity behavior under perturbation
of initial data. The source [E1] verifies the stronger conclusion in
its surface setting; this attempt has no corresponding all-dimensional
stability theorem or replacement for its dimension-specific inputs.

\paragraph{Verification and outcome.}
The local checks verify the sphere and cylinder radius equations, the
Gaussian constants in (3), the exact rational eigenvalues in (5), and
the monotonicity square on evolving round spheres for arbitrary
future Gaussian time. These tests would catch a missing factor of two,
an incorrect curvature sign or the confusion between fixed Gaussian
stability and adjustment of the symmetry parameters. They do not test
weak-flow compactness or genericity.

\textbf{Outcome: UNRESOLVED in the all-dimensional scope.} The surface
case is credited to existing work. The first missing input in the
present approach is an all-dimensional global perturbation and
stability theorem, with multiplicity control and a specified weak-flow
continuation, that promotes local shrinker avoidance to the entire
open dense conclusion. The displayed model and variation calculations
establish only the stated partial checks and reductions.
\subsection{Sources}
\begin{itemize}
\item[S1]T. H. Colding, W. P. Minicozzi II, 'Generic mean curvature flow I; generic singularities', arXiv:0908.3788 (2009); Ann. of Math. (2) 175 (2012) 755-833.
\url{https://arxiv.org/abs/0908.3788}.
\textit{Formulation source.}
\item[E1]Otis Chodosh, Kyeongsu Choi, Christos Mantoulidis and Felix Schulze, \emph{Revisiting generic mean curvature flow in $\mathbb R^3$}, arXiv:2409.01463; Journal f\"ur die reine und angewandte Mathematik 835 (2026), 219--232.
\url{https://arxiv.org/html/2409.01463}.
\textit{Primary text, Theorem 1.1 and Remark 1(c), consulted 27 September 2026 for the resolved surface case, its open dense conclusion and its specified Brakke-flow convention.}
\end{itemize}
\end{document}
